\documentclass[11pt,a4paper]{article}
\usepackage[T1]{fontenc}
\usepackage{isabelle,isabellesym}

\usepackage{amsmath}
\usepackage{amssymb}
\usepackage{stmaryrd}

% this should be the last package used
\usepackage{pdfsetup}

\urlstyle{rm}
\isabellestyle{it}

\begin{document}

\title{Sup-Convolutions, Alexandrov's Theorem and the Crandall--Ishii Theorem on Sums}
\author{Dmitriy Traytel}
\maketitle

\begin{abstract}
  This entry develops the second-order differentiability theory that
  underlies comparison principles for viscosity solutions of fully nonlinear
  degenerate elliptic equations.

  The sup-convolution
  $u^{\varepsilon}(x) = \sup_{y}\,(u(y) - |x-y|^{2}/2\varepsilon)$
  and its convex counterpart, the Moreau envelope, are defined and their
  calculus established: semiconvexity, the proximal map, and the transfer of
  second-order jets between a function and its envelope.  Subgradients of a
  finite convex function on a Euclidean space are shown nonempty by a
  supporting hyperplane to the epigraph, and the proximal map is shown firmly
  nonexpansive.

  On that basis three classical theorems are proved.  \emph{Rademacher's
  theorem}: a locally Lipschitz function on a Euclidean space is
  differentiable almost everywhere, in a scalar and a vector-valued form.
  \emph{Alexandrov's theorem}: a convex function is twice differentiable
  almost everywhere; the proof follows Minty's device, applying Rademacher to
  the resolvent of the subdifferential and transporting the expansion from
  the Moreau envelope back to the function.  \emph{Jensen's lemma} for
  semiconvex functions, and the \emph{Crandall--Ishii theorem on sums}: for a
  local maximum of $u(x)+w(y)$ on a product, second-order jets exist whose
  Hessians satisfy the block matrix inequality.  The development follows
  Evans and Gariepy~\cite{EvansGariepy} for the differentiability theorems
  and the User's Guide of Crandall, Ishii and Lions~\cite{CrandallIshiiLions}
  for Jensen's lemma and the theorem on sums.

  Nothing in the entry refers to any particular equation; it was extracted
  from a formalisation of the uniqueness half of a relative-arbitrage
  problem~\cite{LaiShkolnikovSoner}, where the theorem on sums is what makes
  comparison for a degenerate elliptic operator work.
\end{abstract}

\tableofcontents

\parindent 0pt\parskip 0.5ex

\input{session}

\bibliographystyle{abbrv}
\bibliography{root}

\end{document}
